\documentclass[10pt]{article}
\usepackage{float}
\usepackage{tabularx}
\usepackage{mathtools}
\usepackage{amsmath, amssymb, amsthm}
\usepackage{enumitem}
\usepackage{geometry}
\usepackage{fancyhdr}
\usepackage{titlesec}
\usepackage{tikz}
\usepackage{hyperref}
\usepackage{bookmark}
\geometry{margin=1in}

\pagestyle{fancy}
\fancyhf{}
\title{SFWRENG 3IA3 - Practice Exam}
\author{Matthew Farah, \href{mailto:farahm11@mcmaster.ca}{$\langle$farahm11@mcmaster.ca$\rangle$}, 400468251}
\date{\today}

\hypersetup{
	colorlinks=true,
	linkcolor=blue,
	filecolor=magenta,
	urlcolor=blue,
	citecolor=blue,
	anchorcolor=blue
}

\fancyhead{}
\fancyhead[L]{Matthew Farah, \href{mailto:farahm11@mcmaster.ca}{$\langle$farahm11@mcmaster.ca$\rangle$}, 400468251}
\fancyhead[R]{SFWRENG 3IA3 - Practice Exam}

\setlength{\parindent}{0cm}

\newcommand{\sectiontitle}{SFWRENG 3IA3 - Practice Exam}
\setcounter{section}{1}

\newcounter{question}
\renewcommand{\thequestion}{\arabic{question}}

\newenvironment{question}{
	\newpage
	\refstepcounter{question}
	\phantomsection
	\addcontentsline{toc}{section}{\protect{\large{Question \thequestion}}}
	\vspace{1em}
	\par
	\large\textbf{Question \thequestion}
	\vspace{.2cm}
	\par
	\setcounter{subquestion}{0}
	\setcounter{step}{0}
}{\vspace{.5em}}

\newenvironment{solution}{
	\vspace{1em}\par\large\textbf{{Solution:}}\par\vspace{1em}
}{
	\vspace{1em}
}

\newcounter{subquestion}
\renewcommand{\thesubquestion}{\alph{subquestion}}

\newenvironment{subquestion}{
	\refstepcounter{subquestion}
	\vspace{1em}\par\large\textbf{(\thesubquestion)}\hspace{-.5em}
	\setcounter{step}{0}
	\phantomsection
	\addcontentsline{toc}{subsection}{\protect{\large{\thequestion.\thesubquestion.)}}}
}{
	\par
	\vspace{1em}
}

\makeatother

\makeatletter

\newcounter{step}
\renewcommand{\thestep}{\Roman{step}}

\newenvironment{step}{
	\refstepcounter{step}
	\vspace{1.5em}\par\noindent\large\textbf{(\thestep)}
}{
	\par
	\vspace{1.5em}
}

\makeatother

\newcolumntype{Y}{>{\centering\arraybackslash}X}
\renewcommand{\arraystretch}{1.8}

\fontsize{10pt}{14pt}\selectfont

\begin{document}

\maketitle

\tableofcontents
\newpage

\begin{question}
	Let $X = \left( x_n \right)_{n = 1}^{\infty}$ be a real valued sequence and let $L \in \mathbb{R}$.
\end{question}

\begin{subquestion}
	Define $\lim_{n\to\infty}x_n = L$.
\end{subquestion}

\begin{solution}
	We say that $\lim_{n\to\infty}x_n = L$ if, for arbitrary $\epsilon > 0$, there exists some $N \in \mathbb{N}$ such that for all $n \ge N$, $|x_n - L| < \epsilon$.
\end{solution}

\begin{subquestion}
	Define $X'$ is a subsequence of $X$.
\end{subquestion}

\begin{solution}

\end{solution}

\begin{subquestion}
	Prove that a sequence $X = \left( x_n \right)_{\n = 1}^{\infty}$ converges to $L$ if and only if every subsequence $X'$ of $X$ converges to $L$.
\end{subquestion}

\begin{solution}

\end{solution}

\begin{question}
	Consider the two series $\sum_{n = 1}^{\infty}\cos(n)$ and $\sum_{n = 1}^{\infty}\frac{\cos(n)}{n^2}$.
\end{question}

\begin{subquestion}
	Prove that $\sum_{n = 1}^{\infty}\cos(n)$ is divergent.
\end{subquestion}

\begin{solution}

\end{solution}

\begin{subquestion}
	Prove that $\sum_{n = 1}^{\infty}\frac{\cos(n)}{n^2}$ is convergent.
\end{subquestion}

\begin{solution}

\end{solution}

\begin{question}
	Let $f\colon A \rightarrow \mathbb{R}$ and $c \in A$.
\end{question}

\begin{subquestion}
	Define $f$ is continuous at $c$.
\end{subquestion}

\begin{solution}

\end{solution}

\begin{subquestion}
	Let $f\colon I \rightarrow \mathbb{R}$ be increasing on an interval $I$ and $c$ be an interior point of $I$. Give a rigorous proof that:

	\begin{enumerate}
		\item $\lim_{x\to{c^{-}}}f(x)$ and $\lim_{x\to\c^{+}}f(x)$ both exist.
		\item $\lim_{x\to{c^{-}}}f(x) \le \lim_{x\to{c^{+}}}f(x)$.
		\item $f$ is continuous at $c$ if and only if $\lim_{x\to{c}^{-}}f(x) = f(x) = \lim_{x\to{c}^{+}}f(x)$.
	\end{enumerate}
\end{subquestion}

\begin{solution}

\end{solution}

\begin{question}
	Let $f\colon I \rightarrow \mathbb{R}$ and $c \in I$, where $I$ is an interval.
\end{question}

\begin{subquestion}
	Define differentiability of $f$ at $c$ and state CarathÈodory's Theorem that characterizes differentiability.
\end{subquestion}

\begin{solution}

\end{solution}

\begin{subquestion}
	Suppose in addition that $f$ is strictly increasing and continuous on $I$, that $J \coloneq f(I)$ and that $g\colon J \rightarrow I$ is an inverse function to $f$. If $f$ is differentiable at $c$ and $f'(c) \neq 0$, give a rigorous proof that $g$ is differentiable at $d \coloneq f(c)$ and that:

	\begin{gather*}
		g'(d) = \frac{1}{f'(c)} = \frac{f'(g(d))}.
	\end{gather*}
\end{subquestion}

\begin{solution}

\end{solution}

\begin{question}
	State and prove both the Bolzano-Weierstass Theorem and the Mean Value Theorem.
\end{question}

\begin{solution}

\end{solution}




\end{document}
